\section*{Chapter \ref{ch:nw:where-the-crypto-venues-live} --- Where the Crypto Venues Live}

\begin{solution}{exo:nw:where-the-crypto-venues-live:1}
At most $0.5~\mathrm{ms} \times c_0 / 1.462 = \qty{102.5}{\kilo\metre}$ along the geodesic, and less once equipment and processing are taken out.
\end{solution}

\begin{solution}{exo:nw:where-the-crypto-venues-live:2}
Because the address may belong to the provider's \omterm{def:nw:where-the-crypto-venues-live:as}{edge network} (reached by \omterm{def:nw:where-the-crypto-venues-live:as}{anycast}), which terminates connections near the user and forwards
them to the origin: the lookup then says ``global edge'', not the region of the matching engine.
\end{solution}

\begin{solution}{exo:nw:where-the-crypto-venues-live:3}
The expected best falls from 279 to \qty{185}{\micro\second} over the first ten launches, \qty{94}{\micro\second}; from forty to 160 launches it falls
from 157 to 138, \qty{19}{\micro\second}.
\end{solution}

\begin{solution}{exo:nw:where-the-crypto-venues-live:4}
Zones are kilometres apart, a few microseconds of light, while the uncertainty of a triangulation from thousands of kilometres away (\omterm{def:nw:the-north-american-map:factor}{route factors},
equipment, noise) is tens or hundreds of kilometres. Inside the region, only machines in each zone measuring the venue directly can tell the zones
apart.
\end{solution}

\begin{solution}{exo:nw:where-the-crypto-venues-live:5}
Its round trip to the venue went from within one region to more than \qty{200}{\milli\second} across the world: it could no longer quote competitively.
It should have followed the venue: read the announcement, launched machines in the new region before the cutover, run the lottery there, and moved
its quoting engine.
\end{solution}

\begin{solution}{exo:nw:where-the-crypto-venues-live:6}
It accumulates evidence until the path crosses a boundary set for a 5\% false-alarm rate over the whole sample, so it waits for many days of
evidence; here twelve days for a shift of four standard deviations. Pair it with a fast threshold on the daily or hourly median, confirmed by a
second measurement, and keep the CUSUM test to confirm and to catch slow drifts.
\end{solution}

\begin{solution}{exo:nw:where-the-crypto-venues-live:7}
\texttt{best\_n(LOTTERY, 20.0, 17.14)} gives 23 launches, worth about 1\,830~USD a month net: the optimum moves down and the prize shrinks by a factor of
six.
\end{solution}

\begin{solution}{exo:nw:where-the-crypto-venues-live:8}
The API may be served by an \omterm{def:nw:where-the-crypto-venues-live:as}{edge network} that answers from a \omterm{def:nw:the-european-map:pop}{point of presence} near Frankfurt and forwards to the origin elsewhere; the round trip
measures the distance to the edge, not to the matching engine. Measure what the orders actually traverse (the order-entry path, or the time from an
order to its acknowledgement), and look at the address against the provider's ranges.
\end{solution}

\begin{solution}{pb:nw:where-the-crypto-venues-live:1}
\textbf{Part I.}
\begin{enumerate}
\item \qty{279}{\micro\second} (a lognormal's mean exceeds its median).
\item 185, 157 and \qty{138}{\micro\second}.
\item The lognormal is skewed to the right: rare slow launches pull the mean above the median.
\item Each extra launch saves less than the one before: the expected best falls ever more slowly, so the optimum is finite.
\end{enumerate}
\textbf{Part II.}
\begin{enumerate}[resume]
\item About 11\,400~USD a month.
\item 77 launches, about 11\,700~USD a month net.
\item Very flat: 40 and 120 launches are within a few per cent of the optimum.
\item 23 launches, about 1\,830~USD a month.
\end{enumerate}
\textbf{Part III.}
\begin{enumerate}[resume]
\item Facts: the published median round trip within a zone (\qty{270}{\micro\second}, chapter~16) and the AWS blog's ``sub-150~µs'' intra-Region figure.
  Assumptions: that they are the median and the best 1\% of launches, the lognormal shape, independence, the value and the cost.
\item The provider may place successive launches near each other (or deliberately apart), and capacity changes over the day.
\item The instance type and its network card, its tuning and kernel-bypass stack, the host's other tenants, and the venue's own servers and load
  balancers.
\item Launch, measure each machine's round trip distribution with hardware timestamps, and fit the distribution; repeat on different days.
\end{enumerate}
\textbf{Part IV.}
\begin{enumerate}[resume]
\item \emph{Named result}: with the model's distribution the expected best round trip falls from \qty{279}{\micro\second} for one launch to 185 for ten,
  157 for forty and 138 for 160; at 100~USD a month per microsecond and 17.14~USD a launch, 77 launches maximise the net value, about 11\,700~USD a month.
\item Until the venue or the provider moves something: re-measure every day and restart the lottery on a flag.
\item The zone: the lottery then runs only inside the right zone, and nobody wastes launches in the others.
\item It is open to every user who pays for launches, but it rewards the ones who can afford many; a venue that wants equal access would equalise, as
  chapter~9's venues do.
\item Offer equal-latency access points (private endpoints in each zone with equal paths, chapter~18), or add a delay that removes the difference.
\item Nothing: it launches one machine in the right zone and accepts the median.
\item In the cloud rows: the zone, the machines kept, the re-measurement routine and the launch budget.
\item Because the provider places machines at random distances from the venue, and the best of forty is worth more than the forty launches cost.
\end{enumerate}
\end{solution}

\begin{solution}{iq:nw:where-the-crypto-venues-live:1}
Read the venue's and the provider's public statements; match its endpoint addresses against the provider's published ranges (distinguishing the \omterm{def:nw:where-the-crypto-venues-live:as}{edge
network} from origins); triangulate from machines in several regions; and confirm by launching machines in the candidate region and measuring.
\iqlookfor{Sources with dates; IP ranges and edges; measurement as the final word.}
\end{solution}

\begin{solution}{iq:nw:where-the-crypto-venues-live:2}
One address announced from many places, so that each packet goes to the nearest in the routing system. A round trip to an \omterm{def:nw:where-the-crypto-venues-live:as}{anycast} address measures the
distance to the nearest instance, which may be an edge that forwards to a faraway origin.
\iqlookfor{Routing picks the site; edge against origin; measure the path the orders take.}
\end{solution}

\begin{solution}{iq:nw:where-the-crypto-venues-live:3}
Keep as many as the strategy needs (with a spare), choosing the best by a tail percentile, not a single ping; decide the number of launches by
comparing the expected saving of one more launch with its cost, and re-measure the kept ones every day.
\iqlookfor{Order statistics; tail not mean; cost of search; re-measurement.}
\end{solution}

\begin{solution}{iq:nw:where-the-crypto-venues-live:4}
Track each machine's daily (or hourly) median round trip to the venue; flag a jump beyond a threshold and confirm; run a CUSUM test for slow drifts;
also watch the venue's announcements and the addresses its endpoints resolve to.
\iqlookfor{Fast alarm plus statistical confirmation; announcements; addresses.}
\end{solution}

\begin{solution}{iq:nw:where-the-crypto-venues-live:5}
Asian trading hours carry a large share of crypto volume, several large venues run in AWS's Tokyo region, and firms and other venues followed: a
network effect, which BitMEX cited when it moved from Ireland, and which gives nearby traders a latency edge of hundreds of milliseconds over distant
ones.
\iqlookfor{Network effect; sourced examples; latency asymmetry across regions.}
\end{solution}

\begin{solution}{iq:nw:where-the-crypto-venues-live:6}
Each round trip bounds the distance (half of it at the speed of light in the medium) and, with a \omterm{def:nw:the-north-american-map:factor}{route factor}, estimates it; find the point that best
matches the estimates inside every bound. Accuracy is limited by the unknown \omterm{def:nw:the-north-american-map:factor}{route factors} and equipment delays, by noise, and by the probes'
geometry: from thousands of kilometres away, tens to hundreds of kilometres.
\iqlookfor{Bounds against estimates; \omterm{def:nw:the-north-american-map:factor}{route factor} sensitivity; geometry of probes.}
\end{solution}
